\documentclass[15pt,a4paper]{extarticle}
\usepackage[T2A]{fontenc}
\usepackage[utf8]{inputenc}
\usepackage[russian]{babel}
\usepackage{geometry}
\geometry{left=3cm, right=1.5cm, top=2cm, bottom=2cm}
\usepackage{graphicx}
\usepackage{titlesec}
\usepackage{setspace}
\usepackage{indentfirst}
\usepackage{enumitem}
\usepackage{hyperref}
\usepackage{float}
\usepackage{listings}
\usepackage{xcolor}
\usepackage{ragged2e}
\usepackage{url}  % Добавлено

% Подключаем файл с библиографией


% Убираем рамки и фон для кода
\lstset{
    backgroundcolor=\color{white},
    basicstyle=\ttfamily\footnotesize,
    breaklines=true,
    frame=none,
    tabsize=4,
    xleftmargin=0pt,
    xrightmargin=0pt,
    aboveskip=0pt,
    belowskip=0pt
}

% Убираем рамки вокруг гиперссылок
\hypersetup{
    colorlinks=true,
    linkcolor=black,
    citecolor=black,
    urlcolor=blue,
    allcolors=black,
    pdfborder={0 0 0},
    pdfborderstyle={},
    breaklinks=true,
}

\titleformat{\section}{\normalfont\bfseries\centering}{\thesection}{1em}{}
\titleformat{\subsection}{\normalfont\bfseries}{\thesubsection}{1em}{}
\titleformat{\subsubsection}{\normalfont\itshape}{\thesubsubsection}{1em}{}

% Компактные настройки для списка литературы
\setlist[enumerate]{
    leftmargin=*,
    nosep,
    labelwidth=1.2em,
    leftmargin=2.2em,
    itemindent=0pt,
    parsep=2pt,
    itemsep=3pt
}
\setlength{\emergencystretch}{4em}
\setlength{\parindent}{0pt}
\setlength{\parskip}{3pt}

% Команда для компактного оформления URL
\newcommand{\compacturl}[1]{\url{#1}}

\begin{document}

\begin{titlepage}
    \centering
    \vspace*{1cm}

    \textbf{МИНОБРНАУКИ РОССИИ}
    
    Федеральное государственное бюджетное образовательное учреждение
    
    высшего образования
    
    \textbf{«САРАТОВСКИЙ НАЦИОНАЛЬНЫЙ ИССЛЕДОВАТЕЛЬСКИЙ}
    
    \textbf{ГОСУДАРСТВЕННЫЙ УНИВЕРСИТЕТ}
    
    \textbf{ИМЕНИ Н. Г. ЧЕРНЫШЕВСКОГО»}
    
    \vspace{1.5cm}
    
    \textbf{ПРОЕКТ. ПРИЛОЖЕНИЕ «БЛОГИКУМ».}
    
    \vspace{2cm}
    
    \textbf{ОТЧЁТ}
    
    \vspace{1cm}
    
    \raggedright
    студента 2 курса 211 группы
    
    направления 02.03.02 — Фундаментальная информатика и информационные технологии
    
    факультета КНиИТ
    
    Катушов Егор Степанович
    
    \vspace{2cm}
    
    Проверено:
    
    доцент, к. ф.-м. н. \hfill \underline{\hspace{3cm}} \hfill И. А. Батраева
    
    \vfill
    
    \centering
    Саратов 2025
    
\end{titlepage}

\tableofcontents
\newpage

\section{Цель работы. Постановка задачи}

\textbf{Цель работы:} доработка приложения «Блогикум» по заданию из урока №4 темы №5.

\textbf{Постановка задачи:} необходимо реализовать следующие функциональные возможности:

\begin{enumerate}
    \item Кастомные страницы для ошибок (403, 404, 500).
    \item Работа с пользователями. Подключить к проекту пользователей. Создать страницу \texttt{auth/registration/} с формой для регистрации пользователей. Создать страницу пользователя.
    \item Пагинация. Подключить к проекту пагинацию и настроить вывод не более 10 публикаций: на главную страницу, на страницу пользователя, на страницу категории.
    \item Изображения к постам. Добавить возможность прикреплять изображение к публикациям проекта. Если изображение добавлено, то оно должно отображаться в публикациях на главной странице, странице пользователя, странице категории, отдельной странице публикации.
    \item Добавление новых публикаций. У зарегистрированных пользователей должна быть возможность самостоятельно публиковать посты.
    \item Редактирование публикаций.
    \item Комментарии к публикациям. Создать систему комментирования записей. На странице поста под текстом записи должна выводиться форма для отправки комментария, а ниже — список комментариев.
    \item Удаление публикаций и комментариев. Авторизованные пользователи должны иметь возможность удалять собственные публикации и комментарии. Перед удалением материала должна открываться подтверждающая страница, содержащая публикацию или комментарий.
    \item Новые статичные страницы. Обновить механизм создания и изменения статичных страниц в проекте, используя CBV. Адреса уже существующих статичных страниц не должны измениться.
    \item Отправка электронной почты. Подключить файловый бэкенд: все «отправленные» письма должны аккумулироваться в специальной директории проекта \texttt{sent\_emails}. Из репозитория проекта эту директорию стоит исключить.
    \item Для выполнения задания необходимо использовать view-функции (FBV), view-классы (CBV) или их микс.
\end{enumerate}

\section{Реализация проекта}
\subsection{Кастомные страницы для ошибок}

В проекте реализованы пользовательские страницы ошибок 403 (CSRF), 404 и 500. При возникновении ошибок HTTP 403, 404 или 500 Django вызывает соответствующие обработчики, которые отображают кастомные шаблоны вместо стандартных страниц ошибок. Шаблоны расположены в \texttt{templates/pages/}, обработчики — в \texttt{pages/views.py}.

\begin{lstlisting}[language=Python]
def page_not_found(request, exception):
    return render(request, 'pages/404.html', status=404)

def csrf_failure(request, reason=''):
    return render(request, 'pages/403csrf.html', status=403)

def server_error(request):
    return render(request, 'pages/500.html', status=500)
\end{lstlisting}

Настройки в \texttt{blogicum/urls.py}:

\begin{lstlisting}[language=Python]
handler403 = 'pages.views.permission_denied'
handler404 = 'pages.views.page_not_found'
handler500 = 'pages.views.server_error'
\end{lstlisting}

Настройка CSRF в \texttt{blogicum/settings.py}:

\begin{lstlisting}[language=Python]
CSRF_FAILURE_VIEW = 'pages.views.csrf_failure'
\end{lstlisting}

\newpage
\subsection{Работа с пользователями}

В проекте подключена система аутентификации Django, реализованы страницы регистрации, входа, профиля и редактирования профиля. Используются встроенные механизмы Django для аутентификации. Страница регистрации создает новых пользователей через \texttt{UserCreationForm}. 

\begin{lstlisting}[language=Python]
path('auth/', include('django.contrib.auth.urls')),
path('auth/registration/', CreateView.as_view(
    template_name='registration/registration_form.html',
    form_class=UserCreationForm,
    success_url=reverse_lazy('blog:index'),
), name='registration'),
\end{lstlisting}

\begin{figure}[H]
    \centering
    \includegraphics[width=0.8\textwidth]{Reg.png}
    \caption{Форма регистрации пользователя}
    \label{fig:registration}
\end{figure}

\begin{figure}[H]
    \centering
    \includegraphics[width=0.8\textwidth]{Avto.png}
    \caption{Форма входа в систему}
    \label{fig:login}
\end{figure}

\begin{figure}[H]
    \centering
    \includegraphics[width=0.8\textwidth]{Pol.png}
    \caption{Страница профиля пользователя}
    \label{fig:profile}
\end{figure}

\newpage
\subsection{Пагинация}

В проекте настроена пагинация на главной странице, странице пользователя и странице категории. Объект \texttt{Paginator} разбивает список публикаций на страницы по 10 элементов. Номер текущей страницы передается через GET-параметр \texttt{page}.

\begin{lstlisting}[language=Python]
paginator = Paginator(post_list, 10)
page_number = request.GET.get('page')
page_obj = paginator.get_page(page_number)
\end{lstlisting}

\begin{figure}[H]
    \centering
    \includegraphics[width=0.9\textwidth]{P2.png}
    \label{fig:user_pagination}
\end{figure}

\begin{figure}[H]
    \centering
    \includegraphics[width=0.6\textwidth]{P1.png}
    \caption{Пагинация}
    \label{fig:pagination_nav}
\end{figure}

На фотографиях представлена страница пользователя с пагинацией публикаций. Отображаются заголовки постов, информация об авторе, категории и времени публикации. В нижней части страницы видна панель навигации с элементами "Первая", "<<", "1", "2". На фото показана расширенная панель навигации пагинации с элементами "1", "2", ">>", "Последняя", что позволяет пользователю удобно перемещаться между страницами публикаций.

Пагинация реализована с использованием Django Paginator, который автоматически рассчитывает количество страниц и предоставляет удобный интерфейс для навигации. На каждой странице отображается не более 10 публикаций, что обеспечивает оптимальную производительность и удобство восприятия информации.

\newpage
\subsection{Изображения к постам}

В проекте реализована возможность прикрепления изображений к публикациям. Добавлено поле \texttt{image} в модель \texttt{Post}, настроена работа с медиафайлами. Поле \texttt{ImageField} в модели \texttt{Post} позволяет загружать изображения. Файлы сохраняются в директории \texttt{media/}, настроенной в проекте. Изображения отображаются в шаблонах через тег \texttt{\{\{ post.image.url \}\}}.

\begin{lstlisting}[language=Python]
image = models.ImageField(
    'Изображение',
    upload_to='posts_images',
    blank=True
)
\end{lstlisting}

\begin{lstlisting}[language=Python]
MEDIA_ROOT = BASE_DIR / 'media'
MEDIA_URL = '/media/'
\end{lstlisting}

\begin{figure}[H]
    \centering
    \includegraphics[width=0.8\textwidth]{Post.png}
    \caption{Публикация с изображением}
    \label{fig:post_image}
\end{figure}

\newpage
\subsection{Добавление и редактирование публикаций}

В проекте реализованы создание и редактирование постов с проверкой прав доступа. Классы \texttt{CreateView} и \texttt{UpdateView} обрабатывают создание и редактирование постов. При создании автоматически устанавливается автор, при редактировании проверяется, что пользователь является автором поста.

\begin{lstlisting}[language=Python]
class PostForm(forms.ModelForm):
    class Meta:
        model = Post
        fields = ['title', 'text', 'pub_date', 
                  'image', 'location', 'category', 
                  'is_published']
\end{lstlisting}

\begin{lstlisting}[language=Python]
class PostCreateView(PostsEditMixin, LoginRequiredMixin, CreateView):
    model = Post
    form_class = CreatePostForm

    def form_valid(self, form):
        form.instance.author = self.request.user
        return super().form_valid(form)
\end{lstlisting}

\begin{figure}[H]
    \centering
    \includegraphics[width=0.8\textwidth]{Pub.png}
    \caption{Форма добавления новой публикации}
    \label{fig:add_post}
\end{figure}

\begin{figure}[H]
    \centering
    \includegraphics[width=0.8\textwidth]{Red.png}
    \caption{Форма редактирования существующей публикации}
    \label{fig:edit_post}
\end{figure}

Формы используют Django ModelForm для автоматической генерации полей на основе модели Post. Поле \texttt{pub\_date} поддерживает отложенные публикации — если установить дату в будущем, пост автоматически станет доступен в указанное время.

\newpage
\subsection{Комментарии к публикациям}

В проекте реализована система комментирования с возможностью добавления, редактирования и удаления. Модель \texttt{Comment} связана с \texttt{Post} и \texttt{User} через ForeignKey. Комментарии создаются через форму, которая автоматически связывает комментарий с текущим пользователем и постом.

\begin{lstlisting}[language=Python]
class Comment(models.Model):
    text = models.TextField('Текст комментария')
    post = models.ForeignKey(Post, on_delete=models.CASCADE,
                             related_name='comments')
    author = models.ForeignKey(User, on_delete=models.CASCADE)
    created_at = models.DateTimeField(auto_now_add=True)
\end{lstlisting}

\begin{lstlisting}[language=Python]
class CreateCommentForm(forms.ModelForm):
    class Meta:
        model = Comment
        fields = ("text",)
\end{lstlisting}

\begin{figure}[H]
    \centering
    \includegraphics[width=0.8\textwidth]{Com.png}
    \caption{Страница публикации с формой добавления комментария}
    \label{fig:comments_form}
\end{figure}

Система комментирования обеспечивает интерактивное взаимодействие пользователей с контентом. Форма комментария отображается непосредственно под публикацией, что упрощает процесс добавления комментариев. Реализованная связь между моделями Post и Comment гарантирует корректное отображение комментариев только для соответствующей публикации и обеспечивает целостность данных при удалении поста.

\newpage
\subsection{Удаление публикаций и комментариев}

В проекте реализовано удаление с подтверждением и проверкой прав. \texttt{DeleteView} предоставляет подтверждающую страницу перед удалением. В методе \texttt{delete()} проверяется, что удаляющий является автором объекта, предотвращая несанкционированное удаление.

\begin{lstlisting}[language=Python]
class PostDeleteView(PostsEditMixin, LoginRequiredMixin, DeleteView):
    model = Post
    success_url = reverse_lazy('blog:index')
    pk_url_kwarg = 'post_id'

    def delete(self, request, *args, **kwargs):
        post = get_object_or_404(Post, pk=self.kwargs[self.pk_url_kwarg])
        if self.request.user != post.author:
            return redirect('blog:index')
        return super().delete(request, *args, **kwargs)
\end{lstlisting}

\newpage
\subsection{Новые статичные страницы}

В проекте статичные страницы переведены на CBV. \texttt{TemplateView} отображает указанный шаблон без дополнительной логики. Переход с FBV на CBV улучшает структуру кода и облегчает расширение функциональности.

\begin{lstlisting}[language=Python]
class AboutPage(TemplateView):
    template_name = 'pages/about.html'

class RulesPage(TemplateView):
    template_name = 'pages/rules.html'
\end{lstlisting}

\newpage
\subsection{Отправка электронной почты}

В проекте настроен файловый бэкенд для отправки писем. Файловый бэкенд сохраняет все отправляемые письма в виде файлов в директории \texttt{sent\_emails/}, что позволяет тестировать email-функциональность без настройки реального SMTP-сервера.

\begin{lstlisting}[language=Python]
EMAIL_BACKEND = 'django.core.mail.backends.filebased.EmailBackend'
EMAIL_FILE_PATH = BASE_DIR / 'sent_emails'
\end{lstlisting}

\newpage
\subsection{Миксины для повторного использования кода}

В рамках проекта реализованы миксины для устранения дублирования кода в view-классах. Миксины группируют общие настройки для похожих view-классов. \texttt{PostsEditMixin} определяет общие параметры для работы с постами, \texttt{CommentEditMixin} — для работы с комментариями, включая URL перенаправления.

\begin{lstlisting}[language=Python]
class PostsEditMixin:
    model = Post
    template_name = 'blog/create.html'

class CommentEditMixin:
    model = Comment
    pk_url_kwarg = 'comment_pk'
    template_name = 'blog/comment.html'
    
    def get_success_url(self):
        return reverse('blog:post_detail', args=[self.kwargs['post_id']])
\end{lstlisting}

\newpage
\section{Заключение}

В ходе выполнения проекта были успешно реализованы все поставленные задачи: система пользователей, работа с публикациями и комментариями, пагинация, обработка изображений, кастомные страницы ошибок и отправка электронной почты. Приложение «Блогикум» получило полноценный функционал современной блог-платформы.

Все компоненты были реализованы с использованием современных подходов Django, включая Class-Based Views, миксины для повторного использования кода, и встроенные механизмы аутентификации и пагинации. Проект демонстрирует умение работать с полным циклом разработки веб-приложения на Django.


\newpage
\section*{СПИСОК ИСПОЛЬЗОВАННЫХ ИСТОЧНИКОВ}
\addcontentsline{toc}{section}{СПИСОК ИСПОЛЬЗОВАННЫХ ИСТОЧНИКОВ}

\begin{enumerate}
    \item Прохоренок, Н. А. Python 3: самое необходимое [Текст] / Н. А. Прохоренок. — Санкт-Петербург : БХВ-Петербург, 2016.
    
    \item Сысоева, М. В. Программирование для «нормальных» с нуля на языке Python, Часть 1 [Текст] / М. В. Сысоева, И. В. Сысоев. — Москва : Базальт СПО, МАКС Пресс, 2018.
    
    \item User authentication in Django [Электронный ресурс]. — [Б. м. : б. и.]. — URL: \url{https://docs.djangoproject.com/en/3.2/topics/auth/#installation} (дата обращения: 12.12.2025). Загл. с экр. Яз. англ.
    
    \item Form and field validation [Электронный ресурс]. — [Б. м. : б. и.]. — URL: \url{https://docs.djangoproject.com/en/3.2/ref/forms/validation/} (дата обращения: 16.12.2025). Загл. с экр. Яз. англ.
    
    \item Paginator [Электронный ресурс]. — [Б. м. : б. и.]. — URL: \url{https://docs.djangoproject.com/en/3.2/ref/paginator/#django.core.paginator.Paginator} (дата обращения: 16.12.2025). Загл. с экр. Яз. англ.
    
    \item Урок «Отправка писем. Эмуляция почтового сервера» в теме «Пользователи в Django» на Яндекс.Практикум [Электронный ресурс]. — [Б. м. : б. и.]. — URL: \url{https://practicum.yandex.ru/learn/associated-programs-backend-st-v2/courses/ec89757c-3eb2-4a7c-995d-71561d3b9661/sprints/733852/topics/168d1c84-dc89-4dd9-b18f-7ae332f51a64/lessons/c8931e25-a6ce-472e-bdb1-2ad1bd1f2e1f/} (дата обращения: 13.12.2025). Загл. с экр. Яз. рус.
    
    \item Урок «Настройка страниц входа и выхода пользователей» в теме «Пользователи в Django» на Яндекс.Практикум [Электронный ресурс]. — [Б. м. : б. и.]. — URL: \url{https://practicum.yandex.ru/learn/associated-programs-backend-st-v2/courses/ec89757c-3eb2-4a7c-995d-71561d3b9661/sprints/733852/topics/168d1c84-dc89-4dd9-b18f-7ae332f51a64/lessons/40dfbd65-baa1-42ab-8ebe-122b385cb589/} (дата обращения: 16.12.2025). Загл. с экр. Яз. рус.

\end{enumerate}

\end{document}
\end{document}
